\documentclass[11pt,a4paper]{article}
% ==============================================================
%  Préambule commun aux fiches -- Advanced Probabilities (M1)
%  Usage, dans chaque fiche :
%     \documentclass[11pt,a4paper]{article}
%     % ==============================================================
%  Préambule commun aux fiches -- Advanced Probabilities (M1)
%  Usage, dans chaque fiche :
%     \documentclass[11pt,a4paper]{article}
%     % ==============================================================
%  Préambule commun aux fiches -- Advanced Probabilities (M1)
%  Usage, dans chaque fiche :
%     \documentclass[11pt,a4paper]{article}
%     \input{../preambule.tex}
%     \fiche{C5}{Vecteurs gaussiens}{Gaussian vectors}
% ==============================================================
\usepackage[T1]{fontenc}
\usepackage[utf8]{inputenc}
\usepackage[french]{babel}
\usepackage{lmodern}
\usepackage[a4paper,margin=1.9cm,top=1.7cm,bottom=1.9cm]{geometry}
\usepackage{amsmath,amssymb,amsthm,mathtools}
\usepackage[inline]{enumitem}
\usepackage{booktabs,array,tabularx}
\usepackage{xcolor}
\usepackage[most]{tcolorbox}
\usepackage{fancyhdr}
\usepackage{hyperref}
\hypersetup{colorlinks=true,linkcolor=blue!50!black,urlcolor=blue!50!black}

\setlist{itemsep=1pt,topsep=3pt,parsep=0pt}
\setlength{\parindent}{0pt}
\setlength{\parskip}{3pt}
\emergencystretch=3em

% ---------- Identification de la fiche -------------------------
\newcommand{\ficheid}{}
\newcommand{\fichetitre}{}
\newcommand{\fiche}[3]{%
  \renewcommand{\ficheid}{#1}\renewcommand{\fichetitre}{#2}%
  \pagestyle{fancy}\fancyhf{}%
  \fancyhead[L]{\small\textsf{Advanced Probabilities -- Fiche #1}}%
  \fancyhead[R]{\small\textsf{#2}}%
  \fancyfoot[C]{\small\thepage}%
  \begin{center}
    {\LARGE\bfseries Fiche #1 -- #2}\par\smallskip
    {\large\itshape (#3)}\par\smallskip
    {\small Cours de V.\,M.~Lambert, M1 -- semestre 1}
  \end{center}\vspace{-2pt}\hrule\vspace{6pt}}

% ---------- Environnements numérotés (un seul compteur) --------
\theoremstyle{definition}
\newtheorem{definition}{Définition}
\renewcommand{\thedefinition}{\ficheid.\arabic{definition}}
\newtheorem{theoreme}[definition]{Théorème}
\newtheorem{proposition}[definition]{Proposition}
\newtheorem{lemme}[definition]{Lemme}
\newtheorem{corollaire}[definition]{Corollaire}
\newtheorem{exemple}[definition]{Exemple}
\theoremstyle{remark}
\newtheorem{remarque}[definition]{Remarque}

% Résultat hors cours (standard, signalé comme tel)
\theoremstyle{definition}
\newtheorem{horscours}[definition]{Hors cours}

% Idée de preuve
\newenvironment{idee}{\par\small\textit{Idée de preuve.}~}{\par}

% ---------- Encadrés -----------------------------------------
\tcbset{fichebox/.style={enhanced,breakable,boxrule=0.6pt,arc=2pt,
  left=5pt,right=5pt,top=3pt,bottom=3pt,fonttitle=\bfseries}}
\newtcolorbox{piege}[1][]{fichebox,colback=red!4,colframe=red!60!black,
  title={Pièges classiques},#1}
\newtcolorbox{liens}[1][]{fichebox,colback=blue!3,colframe=blue!50!black,
  title={Liens entre notions},#1}
\newtcolorbox{methode}[1][Méthode]{fichebox,colback=green!4,colframe=green!40!black,title={#1}}
\newtcolorbox{alerte}[1][]{fichebox,colback=orange!6,colframe=orange!70!black,
  title={Attention à l'énoncé},#1}

% Exercice de TD (fiches de rappels)
\newcounter{exo}
\newenvironment{exo}[1][]{\refstepcounter{exo}\par\medskip
  \noindent{\large\bfseries Exercice~\theexo}\ifx\relax#1\relax\else~{\large\bfseries -- #1}\fi\par\nopagebreak\smallskip}{\par}
\newcommand{\outils}{\textbf{Outils du cours.}~}
\newcommand{\meth}{\textbf{Méthode / astuce.}~}
\newcommand{\indic}{\textbf{Indication.}~}
\newcommand{\rep}{\textbf{Réponse.}~}

% ---------- Notations du cours ---------------------------------
\newcommand{\R}{\mathbb{R}}
\newcommand{\N}{\mathbb{N}}
\newcommand{\Q}{\mathbb{Q}}
\newcommand{\C}{\mathbb{C}}
\newcommand{\PP}{\mathbb{P}}
\newcommand{\QQ}{\mathbb{Q}}
\newcommand{\E}{\mathbb{E}}
\newcommand{\F}{\mathcal{F}}
\newcommand{\G}{\mathcal{G}}
\newcommand{\Hc}{\mathcal{H}}
\newcommand{\Tc}{\mathcal{T}}
\newcommand{\Bc}{\mathcal{B}}
\newcommand{\Cc}{\mathcal{C}}
\newcommand{\Ec}{\mathcal{E}}
\newcommand{\Pc}{\mathcal{P}}
\newcommand{\Nc}{\mathcal{N}}
\newcommand{\LN}{\mathcal{LN}}
\newcommand{\Expo}{\mathcal{E}}
\newcommand{\Poi}{\mathcal{P}}
\newcommand{\Ff}{\mathbb{F}}
\newcommand{\1}{\mathbf{1}}
\newcommand{\dd}{\,\mathrm{d}}
\DeclareMathOperator{\Var}{\mathbb{V}}
\DeclareMathOperator{\Cov}{Cov}
\DeclareMathOperator{\card}{card}
\DeclareMathOperator{\diag}{diag}
\newcommand{\indep}{\perp\!\!\!\perp}
\newcommand{\ps}{\ \PP\text{-p.s.}}
\newcommand{\tendloi}{\xrightarrow[n\to+\infty]{\text{loi}}}
\newcommand{\tendps}{\xrightarrow[n\to+\infty]{\text{p.s.}}}
\newcommand{\tendproba}{\xrightarrow[n\to+\infty]{\PP}}
\newcommand{\tendLp}{\xrightarrow[n\to+\infty]{L^p}}
\newcommand{\tend}{\xrightarrow[n\to+\infty]{}}

%     \fiche{C5}{Vecteurs gaussiens}{Gaussian vectors}
% ==============================================================
\usepackage[T1]{fontenc}
\usepackage[utf8]{inputenc}
\usepackage[french]{babel}
\usepackage{lmodern}
\usepackage[a4paper,margin=1.9cm,top=1.7cm,bottom=1.9cm]{geometry}
\usepackage{amsmath,amssymb,amsthm,mathtools}
\usepackage[inline]{enumitem}
\usepackage{booktabs,array,tabularx}
\usepackage{xcolor}
\usepackage[most]{tcolorbox}
\usepackage{fancyhdr}
\usepackage{hyperref}
\hypersetup{colorlinks=true,linkcolor=blue!50!black,urlcolor=blue!50!black}

\setlist{itemsep=1pt,topsep=3pt,parsep=0pt}
\setlength{\parindent}{0pt}
\setlength{\parskip}{3pt}
\emergencystretch=3em

% ---------- Identification de la fiche -------------------------
\newcommand{\ficheid}{}
\newcommand{\fichetitre}{}
\newcommand{\fiche}[3]{%
  \renewcommand{\ficheid}{#1}\renewcommand{\fichetitre}{#2}%
  \pagestyle{fancy}\fancyhf{}%
  \fancyhead[L]{\small\textsf{Advanced Probabilities -- Fiche #1}}%
  \fancyhead[R]{\small\textsf{#2}}%
  \fancyfoot[C]{\small\thepage}%
  \begin{center}
    {\LARGE\bfseries Fiche #1 -- #2}\par\smallskip
    {\large\itshape (#3)}\par\smallskip
    {\small Cours de V.\,M.~Lambert, M1 -- semestre 1}
  \end{center}\vspace{-2pt}\hrule\vspace{6pt}}

% ---------- Environnements numérotés (un seul compteur) --------
\theoremstyle{definition}
\newtheorem{definition}{Définition}
\renewcommand{\thedefinition}{\ficheid.\arabic{definition}}
\newtheorem{theoreme}[definition]{Théorème}
\newtheorem{proposition}[definition]{Proposition}
\newtheorem{lemme}[definition]{Lemme}
\newtheorem{corollaire}[definition]{Corollaire}
\newtheorem{exemple}[definition]{Exemple}
\theoremstyle{remark}
\newtheorem{remarque}[definition]{Remarque}

% Résultat hors cours (standard, signalé comme tel)
\theoremstyle{definition}
\newtheorem{horscours}[definition]{Hors cours}

% Idée de preuve
\newenvironment{idee}{\par\small\textit{Idée de preuve.}~}{\par}

% ---------- Encadrés -----------------------------------------
\tcbset{fichebox/.style={enhanced,breakable,boxrule=0.6pt,arc=2pt,
  left=5pt,right=5pt,top=3pt,bottom=3pt,fonttitle=\bfseries}}
\newtcolorbox{piege}[1][]{fichebox,colback=red!4,colframe=red!60!black,
  title={Pièges classiques},#1}
\newtcolorbox{liens}[1][]{fichebox,colback=blue!3,colframe=blue!50!black,
  title={Liens entre notions},#1}
\newtcolorbox{methode}[1][Méthode]{fichebox,colback=green!4,colframe=green!40!black,title={#1}}
\newtcolorbox{alerte}[1][]{fichebox,colback=orange!6,colframe=orange!70!black,
  title={Attention à l'énoncé},#1}

% Exercice de TD (fiches de rappels)
\newcounter{exo}
\newenvironment{exo}[1][]{\refstepcounter{exo}\par\medskip
  \noindent{\large\bfseries Exercice~\theexo}\ifx\relax#1\relax\else~{\large\bfseries -- #1}\fi\par\nopagebreak\smallskip}{\par}
\newcommand{\outils}{\textbf{Outils du cours.}~}
\newcommand{\meth}{\textbf{Méthode / astuce.}~}
\newcommand{\indic}{\textbf{Indication.}~}
\newcommand{\rep}{\textbf{Réponse.}~}

% ---------- Notations du cours ---------------------------------
\newcommand{\R}{\mathbb{R}}
\newcommand{\N}{\mathbb{N}}
\newcommand{\Q}{\mathbb{Q}}
\newcommand{\C}{\mathbb{C}}
\newcommand{\PP}{\mathbb{P}}
\newcommand{\QQ}{\mathbb{Q}}
\newcommand{\E}{\mathbb{E}}
\newcommand{\F}{\mathcal{F}}
\newcommand{\G}{\mathcal{G}}
\newcommand{\Hc}{\mathcal{H}}
\newcommand{\Tc}{\mathcal{T}}
\newcommand{\Bc}{\mathcal{B}}
\newcommand{\Cc}{\mathcal{C}}
\newcommand{\Ec}{\mathcal{E}}
\newcommand{\Pc}{\mathcal{P}}
\newcommand{\Nc}{\mathcal{N}}
\newcommand{\LN}{\mathcal{LN}}
\newcommand{\Expo}{\mathcal{E}}
\newcommand{\Poi}{\mathcal{P}}
\newcommand{\Ff}{\mathbb{F}}
\newcommand{\1}{\mathbf{1}}
\newcommand{\dd}{\,\mathrm{d}}
\DeclareMathOperator{\Var}{\mathbb{V}}
\DeclareMathOperator{\Cov}{Cov}
\DeclareMathOperator{\card}{card}
\DeclareMathOperator{\diag}{diag}
\newcommand{\indep}{\perp\!\!\!\perp}
\newcommand{\ps}{\ \PP\text{-p.s.}}
\newcommand{\tendloi}{\xrightarrow[n\to+\infty]{\text{loi}}}
\newcommand{\tendps}{\xrightarrow[n\to+\infty]{\text{p.s.}}}
\newcommand{\tendproba}{\xrightarrow[n\to+\infty]{\PP}}
\newcommand{\tendLp}{\xrightarrow[n\to+\infty]{L^p}}
\newcommand{\tend}{\xrightarrow[n\to+\infty]{}}

%     \fiche{C5}{Vecteurs gaussiens}{Gaussian vectors}
% ==============================================================
\usepackage[T1]{fontenc}
\usepackage[utf8]{inputenc}
\usepackage[french]{babel}
\usepackage{lmodern}
\usepackage[a4paper,margin=1.9cm,top=1.7cm,bottom=1.9cm]{geometry}
\usepackage{amsmath,amssymb,amsthm,mathtools}
\usepackage[inline]{enumitem}
\usepackage{booktabs,array,tabularx}
\usepackage{xcolor}
\usepackage[most]{tcolorbox}
\usepackage{fancyhdr}
\usepackage{hyperref}
\hypersetup{colorlinks=true,linkcolor=blue!50!black,urlcolor=blue!50!black}

\setlist{itemsep=1pt,topsep=3pt,parsep=0pt}
\setlength{\parindent}{0pt}
\setlength{\parskip}{3pt}
\emergencystretch=3em

% ---------- Identification de la fiche -------------------------
\newcommand{\ficheid}{}
\newcommand{\fichetitre}{}
\newcommand{\fiche}[3]{%
  \renewcommand{\ficheid}{#1}\renewcommand{\fichetitre}{#2}%
  \pagestyle{fancy}\fancyhf{}%
  \fancyhead[L]{\small\textsf{Advanced Probabilities -- Fiche #1}}%
  \fancyhead[R]{\small\textsf{#2}}%
  \fancyfoot[C]{\small\thepage}%
  \begin{center}
    {\LARGE\bfseries Fiche #1 -- #2}\par\smallskip
    {\large\itshape (#3)}\par\smallskip
    {\small Cours de V.\,M.~Lambert, M1 -- semestre 1}
  \end{center}\vspace{-2pt}\hrule\vspace{6pt}}

% ---------- Environnements numérotés (un seul compteur) --------
\theoremstyle{definition}
\newtheorem{definition}{Définition}
\renewcommand{\thedefinition}{\ficheid.\arabic{definition}}
\newtheorem{theoreme}[definition]{Théorème}
\newtheorem{proposition}[definition]{Proposition}
\newtheorem{lemme}[definition]{Lemme}
\newtheorem{corollaire}[definition]{Corollaire}
\newtheorem{exemple}[definition]{Exemple}
\theoremstyle{remark}
\newtheorem{remarque}[definition]{Remarque}

% Résultat hors cours (standard, signalé comme tel)
\theoremstyle{definition}
\newtheorem{horscours}[definition]{Hors cours}

% Idée de preuve
\newenvironment{idee}{\par\small\textit{Idée de preuve.}~}{\par}

% ---------- Encadrés -----------------------------------------
\tcbset{fichebox/.style={enhanced,breakable,boxrule=0.6pt,arc=2pt,
  left=5pt,right=5pt,top=3pt,bottom=3pt,fonttitle=\bfseries}}
\newtcolorbox{piege}[1][]{fichebox,colback=red!4,colframe=red!60!black,
  title={Pièges classiques},#1}
\newtcolorbox{liens}[1][]{fichebox,colback=blue!3,colframe=blue!50!black,
  title={Liens entre notions},#1}
\newtcolorbox{methode}[1][Méthode]{fichebox,colback=green!4,colframe=green!40!black,title={#1}}
\newtcolorbox{alerte}[1][]{fichebox,colback=orange!6,colframe=orange!70!black,
  title={Attention à l'énoncé},#1}

% Exercice de TD (fiches de rappels)
\newcounter{exo}
\newenvironment{exo}[1][]{\refstepcounter{exo}\par\medskip
  \noindent{\large\bfseries Exercice~\theexo}\ifx\relax#1\relax\else~{\large\bfseries -- #1}\fi\par\nopagebreak\smallskip}{\par}
\newcommand{\outils}{\textbf{Outils du cours.}~}
\newcommand{\meth}{\textbf{Méthode / astuce.}~}
\newcommand{\indic}{\textbf{Indication.}~}
\newcommand{\rep}{\textbf{Réponse.}~}

% ---------- Notations du cours ---------------------------------
\newcommand{\R}{\mathbb{R}}
\newcommand{\N}{\mathbb{N}}
\newcommand{\Q}{\mathbb{Q}}
\newcommand{\C}{\mathbb{C}}
\newcommand{\PP}{\mathbb{P}}
\newcommand{\QQ}{\mathbb{Q}}
\newcommand{\E}{\mathbb{E}}
\newcommand{\F}{\mathcal{F}}
\newcommand{\G}{\mathcal{G}}
\newcommand{\Hc}{\mathcal{H}}
\newcommand{\Tc}{\mathcal{T}}
\newcommand{\Bc}{\mathcal{B}}
\newcommand{\Cc}{\mathcal{C}}
\newcommand{\Ec}{\mathcal{E}}
\newcommand{\Pc}{\mathcal{P}}
\newcommand{\Nc}{\mathcal{N}}
\newcommand{\LN}{\mathcal{LN}}
\newcommand{\Expo}{\mathcal{E}}
\newcommand{\Poi}{\mathcal{P}}
\newcommand{\Ff}{\mathbb{F}}
\newcommand{\1}{\mathbf{1}}
\newcommand{\dd}{\,\mathrm{d}}
\DeclareMathOperator{\Var}{\mathbb{V}}
\DeclareMathOperator{\Cov}{Cov}
\DeclareMathOperator{\card}{card}
\DeclareMathOperator{\diag}{diag}
\newcommand{\indep}{\perp\!\!\!\perp}
\newcommand{\ps}{\ \PP\text{-p.s.}}
\newcommand{\tendloi}{\xrightarrow[n\to+\infty]{\text{loi}}}
\newcommand{\tendps}{\xrightarrow[n\to+\infty]{\text{p.s.}}}
\newcommand{\tendproba}{\xrightarrow[n\to+\infty]{\PP}}
\newcommand{\tendLp}{\xrightarrow[n\to+\infty]{L^p}}
\newcommand{\tend}{\xrightarrow[n\to+\infty]{}}

% Encadré local pour les résultats (macro absente du préambule)
\newtcolorbox{resultat}{fichebox,colback=gray!6,colframe=black!55}
\newcommand{\question}[1]{\par\medskip\noindent\textbf{Question #1.}~}
\begin{document}
\fiche{TD1}{Corrigé -- Tribus et mesures de probabilité}{Sigma-algebra, Probability Measure}

\textit{Corrigé rédigé du TD1. Renvois : « C1.4 » = fiche de cours C1, énoncé n°4 ; « C1.11, point 2 » = item 2 de
la proposition C1.11. Tout au long, $\Omega$ est un ensemble non vide et « tribu » = \emph{sigma-algebra} au sens de C1.1 :
contient $\emptyset$, stable par complémentaire (axiome (ii)) et par union dénombrable (axiome (iii)).}

% =====================================================================
\begin{exo}[Tribu engendrée par un ensemble]

\question{1} \emph{Soit $(\Tc_i)_{i\in I}$ une famille de tribus sur $\Omega$. Montrer que $\Tc:=\bigcap_{i\in I}\Tc_i$ est une tribu.}

On suppose $I\neq\emptyset$ ; $I$ est \emph{quelconque} (fini, dénombrable ou non). (Pour $I=\emptyset$, la convention
$\bigcap_{i\in\emptyset}\Tc_i=\Pc(\Omega)$ donne encore une tribu, C1.3, point 2.) Rappelons que
$A\in\Tc\iff\forall i\in I,\ A\in\Tc_i$. On vérifie les trois axiomes de C1.1.
\begin{enumerate}[label=(\roman*)]
  \item Pour tout $i$, $\emptyset\in\Tc_i$ car $\Tc_i$ est une tribu (C1.1 (i)) ; donc $\emptyset\in\Tc$.
  \item Soit $A\in\Tc$. Pour tout $i$, $A\in\Tc_i$, donc $A^c\in\Tc_i$ (C1.1 (ii) pour $\Tc_i$). Ainsi $A^c\in\Tc_i$ pour tout $i$, i.e.\ $A^c\in\Tc$.
  \item Soit $(A_n)_{n\in\N}\subset\Tc$. Pour $i$ fixé, $(A_n)_n\subset\Tc_i$, donc $\bigcup_nA_n\in\Tc_i$ (C1.1 (iii) pour $\Tc_i$).
  C'est vrai pour tout $i$, donc $\bigcup_nA_n\in\Tc$.
\end{enumerate}
Enfin $\Tc\subset\Tc_{i_0}\subset\Pc(\Omega)$ pour un $i_0\in I$.
\begin{resultat}
Toute intersection (même non dénombrable) de tribus sur $\Omega$ est une tribu sur $\Omega$.
\end{resultat}

\question{2} \emph{Exemple de deux tribus $\Tc_1,\Tc_2$ telles que $\Tc_1\cup\Tc_2$ ne soit pas une tribu.}

Soit $\Omega=\{1,2,3\}$, $\Tc_1=\{\emptyset,\{1\},\{2,3\},\Omega\}$ et $\Tc_2=\{\emptyset,\{2\},\{1,3\},\Omega\}$.
Ce sont des tribus : ce sont les tribus des partitions $\{\{1\},\{2,3\}\}$ et $\{\{2\},\{1,3\}\}$ (C1.3, point 3, et
exercice 2 ci-dessous), c'est-à-dire $\Tc_1=\sigma(\{\{1\}\})$, $\Tc_2=\sigma(\{\{2\}\})$. Alors
\[\Tc_1\cup\Tc_2=\bigl\{\emptyset,\{1\},\{2\},\{2,3\},\{1,3\},\Omega\bigr\}.\]
On a $\{1\}\in\Tc_1\cup\Tc_2$ et $\{2\}\in\Tc_1\cup\Tc_2$, mais $\{1\}\cup\{2\}=\{1,2\}\notin\Tc_1\cup\Tc_2$.
Or une tribu est stable par union finie (remarque suivant C1.1 : compléter la suite par des $\emptyset$ et appliquer (iii)).
Donc $\Tc_1\cup\Tc_2$ n'est pas une tribu. (Elle contient $\emptyset$ et est stable par complémentaire : seul l'axiome (iii) échoue.)
\begin{resultat}
$\Omega=\{1,2,3\}$, $\Tc_1=\sigma(\{1\})$, $\Tc_2=\sigma(\{2\})$ : $\{1,2\}\notin\Tc_1\cup\Tc_2$, donc $\Tc_1\cup\Tc_2$ n'est pas une tribu.
\end{resultat}

% À VÉRIFIER : interprétation de T1 x T2 comme famille des pavés.
\question{3} \emph{Exemple de deux tribus $\Tc_1,\Tc_2$ telles que $\Tc_1\times\Tc_2$ ne soit pas une tribu.}

\begin{alerte}
L'énoncé ne précise pas les ensembles. On prend $\Tc_1$ tribu sur $\Omega_1$, $\Tc_2$ tribu sur $\Omega_2$, et on interprète
$\Tc_1\times\Tc_2$ comme la famille des \emph{pavés} $\{A\times B,\ A\in\Tc_1,\ B\in\Tc_2\}\subset\Pc(\Omega_1\times\Omega_2)$
(le produit cartésien des familles, ensemble de couples $(A,B)$, n'est pas une famille de parties d'un ensemble, donc la question
n'aurait pas de sens). C'est l'interprétation de C3.12.
\end{alerte}

Soit $\Omega_1=\Omega_2=\{0,1\}$ et $\Tc_1=\Tc_2=\Pc(\{0,1\})$ (tribus par C1.3, point 2). Le pavé $C=\{0\}\times\{0\}=\{(0,0)\}$
appartient à $\Tc_1\times\Tc_2$. Son complémentaire dans $\Omega_1\times\Omega_2$ est
\[C^c=\{(0,1),(1,0),(1,1)\}.\]
Supposons $C^c=A\times B$ avec $A,B\subset\{0,1\}$. Comme $(0,1)\in A\times B$, on a $0\in A$ ; comme $(1,0)\in A\times B$, on a $0\in B$.
Donc $(0,0)\in A\times B=C^c$, contradiction. Ainsi $C^c\notin\Tc_1\times\Tc_2$ : l'axiome (ii) de C1.1 est violé.
\begin{resultat}
Avec $\Tc_1=\Tc_2=\Pc(\{0,1\})$, le complémentaire de $\{(0,0)\}$ n'est pas un pavé : $\Tc_1\times\Tc_2$ n'est pas une tribu.
C'est pourquoi la tribu produit est \emph{définie} par $\Tc_1\otimes\Tc_2:=\sigma(\Tc_1\times\Tc_2)$ (C3.12).
\end{resultat}
\textit{Remarque.} $\Tc_1\times\Tc_2$ est tout de même stable par intersection finie : $(A\times B)\cap(A'\times B')=(A\cap A')\times(B\cap B')$.

\question{4} \emph{Pour $\Cc\subset\Pc(\Omega)$, montrer qu'il existe une plus petite tribu contenant $\Cc$.}

Soit $\mathcal S=\{\Tc\ \text{tribu sur }\Omega,\ \Cc\subset\Tc\}$.
\begin{itemize}
  \item $\mathcal S\neq\emptyset$ : $\Pc(\Omega)$ est une tribu (C1.3, point 2) et contient $\Cc$.
  \item Posons $\sigma(\Cc):=\bigcap_{\Tc\in\mathcal S}\Tc$ (famille de tribus indexée par $I=\mathcal S\neq\emptyset$).
  Par la question 1, $\sigma(\Cc)$ est une tribu.
  \item Elle contient $\Cc$ : si $C\in\Cc$, alors $C\in\Tc$ pour tout $\Tc\in\mathcal S$, donc $C\in\sigma(\Cc)$. Ainsi $\sigma(\Cc)\in\mathcal S$.
  \item Elle est la plus petite : si $\Tc'$ est une tribu contenant $\Cc$, alors $\Tc'\in\mathcal S$, donc $\sigma(\Cc)=\bigcap_{\Tc\in\mathcal S}\Tc\subset\Tc'$.
  \item Unicité : si $\Tc^*$ et $\Tc^{**}$ sont deux plus petits éléments de $\mathcal S$, on a $\Tc^*\subset\Tc^{**}$ et $\Tc^{**}\subset\Tc^*$, donc $\Tc^*=\Tc^{**}$.
\end{itemize}
\begin{resultat}
$\sigma(\Cc)=\bigcap\{\Tc\ \text{tribu},\ \Cc\subset\Tc\}$ est l'unique plus petite tribu contenant $\Cc$ (C1.4).
Conséquence utile : si $\Tc$ est une tribu et $\Cc\subset\Tc$, alors $\sigma(\Cc)\subset\Tc$.
\end{resultat}

\question{5} \emph{Contre-exemple à $\sigma(\Cc_1)\cap\sigma(\Cc_2)=\sigma(\Cc_1\cap\Cc_2)$.}

Soit $\Omega=\{1,2\}$ (ou tout $\Omega$ avec $A\notin\{\emptyset,\Omega\}$), $A=\{1\}$, $\Cc_1=\{A\}$ et $\Cc_2=\{A^c\}=\{\{2\}\}$.
\begin{itemize}
  \item $\sigma(\Cc_1)=\{\emptyset,A,A^c,\Omega\}$ : cette famille est une tribu (tribu de la partition $\{A,A^c\}$, C1.3, point 3),
  elle contient $A$, et toute tribu contenant $A$ contient $A^c$ (axiome (ii)), $\emptyset$ (axiome (i)) et $\Omega=\emptyset^c$ ;
  on conclut par la minimalité de C1.4. De même $\sigma(\Cc_2)=\{\emptyset,A^c,A,\Omega\}$. Donc
  $\sigma(\Cc_1)\cap\sigma(\Cc_2)=\{\emptyset,A,A^c,\Omega\}$ ($=\Pc(\Omega)$ si $\Omega=\{1,2\}$).
  \item $\Cc_1\cap\Cc_2=\emptyset$ (famille vide) car $A\neq A^c$. Or $\sigma(\emptyset)=\{\emptyset,\Omega\}$ :
  la tribu triviale (C1.3, point 1) contient la famille vide, et toute tribu contient $\emptyset$ et $\Omega$, donc la contient.
\end{itemize}
Comme $A\notin\{\emptyset,\Omega\}$, les deux tribus diffèrent.
\begin{resultat}
$\Cc_1=\{A\}$, $\Cc_2=\{A^c\}$ avec $A\neq\emptyset,\Omega$ :
$\sigma(\Cc_1)\cap\sigma(\Cc_2)=\{\emptyset,A,A^c,\Omega\}\neq\{\emptyset,\Omega\}=\sigma(\Cc_1\cap\Cc_2)$.
\end{resultat}
\textit{Remarque.} L'inclusion $\sigma(\Cc_1\cap\Cc_2)\subset\sigma(\Cc_1)\cap\sigma(\Cc_2)$ est, elle, toujours vraie :
$\sigma(\Cc_1)\cap\sigma(\Cc_2)$ est une tribu (question 1) contenant $\Cc_1\cap\Cc_2$, puis minimalité (C1.4).
\end{exo}

% =====================================================================
\begin{exo}[Tribu engendrée par une partition]
Soit $(A_i)_{i\in I}$ une partition dénombrable de $\Omega$ : $I$ est au plus dénombrable, les $A_i$ sont deux à deux disjoints
et $\bigcup_{i\in I}A_i=\Omega$. Ainsi tout $\omega\in\Omega$ appartient à \emph{exactement un} $A_i$ ; notons $i(\omega)$ cet indice.
On note $A_J:=\bigcup_{j\in J}A_j$ pour $J\subset I$, de sorte que $\Tc=\{A_J,\ J\subset I\}$.

\question{1} \emph{Montrer que $\Tc$ est une tribu.}

Observation clé : pour $J\subset I$, $\omega\in A_J\iff i(\omega)\in J$ (si $\omega\in A_j$ avec $j\in J$, alors $j=i(\omega)$ par unicité).
\begin{enumerate}[label=(\roman*)]
  \item $J=\emptyset$ donne $A_\emptyset=\emptyset$ (union vide), donc $\emptyset\in\Tc$.
  \item Pour $J\subset I$ : $\omega\in(A_J)^c\iff i(\omega)\notin J\iff i(\omega)\in I\setminus J\iff\omega\in A_{I\setminus J}$.
  Donc $(A_J)^c=A_{I\setminus J}\in\Tc$.
  \item Pour $(J_n)_{n\in\N}$ parties de $I$ : $\bigcup_nA_{J_n}=\bigcup_n\bigcup_{j\in J_n}A_j=\bigcup_{j\in J}A_j=A_J$ avec
  $J:=\bigcup_nJ_n\subset I$. Donc $\bigcup_nA_{J_n}\in\Tc$.
\end{enumerate}
\begin{resultat}
$\Tc=\bigl\{\bigcup_{j\in J}A_j,\ J\subset I\bigr\}$ est une tribu (c'est la tribu de partition de C1.3, point 3).
\end{resultat}
\textit{Remarque.} Cette question n'utilise pas la dénombrabilité de $I$ ; seule l'hypothèse de \emph{partition} sert (au point (ii)).

\question{2} \emph{Montrer que $\Tc=\sigma(\{A_i,\ i\in I\})$.}

Par double inclusion.
\begin{itemize}
  \item $\sigma(\{A_i\})\subset\Tc$ : $\Tc$ est une tribu (question 1) qui contient chaque $A_i=A_{\{i\}}$. Par minimalité de la
  tribu engendrée (C1.4, ex.~1 question 4), $\sigma(\{A_i,\ i\in I\})\subset\Tc$.
  \item $\Tc\subset\sigma(\{A_i\})$ : soit $J\subset I$. Comme $I$ est au plus dénombrable, $J$ aussi. Si $J=\emptyset$, $A_J=\emptyset\in\sigma(\{A_i\})$.
  Sinon, on énumère $J=\{j_0,j_1,\dots\}$ (en répétant un indice si $J$ est fini) ; $A_J=\bigcup_{n\in\N}A_{j_n}$ est une union
  dénombrable d'éléments de $\sigma(\{A_i\})$, donc appartient à $\sigma(\{A_i\})$ (C1.1 (iii)).
\end{itemize}
\begin{resultat}
$\Tc=\sigma(\{A_i,\ i\in I\})$ : les événements de la tribu engendrée par une partition dénombrable sont exactement les unions de blocs.
\end{resultat}
\textit{Remarque.} La dénombrabilité de $I$ est indispensable pour la seconde inclusion : pour la partition de $\R$ en singletons,
$\sigma(\{\{x\},\ x\in\R\})$ est la tribu des parties dénombrables ou de complémentaire dénombrable, strictement plus petite que $\Tc=\Pc(\R)$
(par exemple $[0,1]\notin\sigma(\{\{x\}\})$). \emph{Hors cours} : ce fait n'est pas démontré dans le cours.
\end{exo}

% =====================================================================
\begin{exo}[Compatibilité de probabilités]
\question{a} \emph{Peut-on construire une probabilité $\PP$ sur $(\Omega,\Tc)$ avec $\PP(A)=0{,}7$, $\PP(B)=0{,}9$, $\PP(A\cap B)=0{,}5$ ?}

On a $A,B\in\Tc$, donc $A\cap B$ et $A\cup B$ sont dans $\Tc$ (remarque suivant C1.1), et $\PP(A\cup B)$ est défini.
Supposons qu'une telle $\PP$ existe. Par C1.11, point 2 (formule d'inclusion--exclusion, conséquence de la $\sigma$-additivité appliquée à
$A\cup B=A\sqcup(B\setminus A)$ et $B=(A\cap B)\sqcup(B\setminus A)$) :
\[\PP(A\cup B)=\PP(A)+\PP(B)-\PP(A\cap B)=0{,}7+0{,}9-0{,}5=1{,}1.\]
Or $A\cup B\subset\Omega$, donc par monotonie (C1.11, point 1) et C1.10 (i), $\PP(A\cup B)\le\PP(\Omega)=1$. Contradiction.
(Les autres contraintes sont satisfaites : $\PP(A\cap B)=0{,}5\le\min(0{,}7;0{,}9)$ ; c'est uniquement la borne sur l'union qui échoue.)
\begin{resultat}
Impossible, quels que soient $\Omega$, $\Tc$, $A$, $B$ : on aurait $\PP(A\cup B)=1{,}1>1$.
De façon équivalente, C1.11, point 2 impose $\PP(A\cap B)\ge\PP(A)+\PP(B)-1=0{,}6>0{,}5$.
\end{resultat}

\question{b} \emph{Minorer $\PP(A\cup B)$ à l'aide de $\PP(A)$ et $\PP(B)$, pour tous $A,B\in\Tc$.}

Comme $A\subset A\cup B$ et $B\subset A\cup B$, la monotonie (C1.11, point 1) donne $\PP(A\cup B)\ge\PP(A)$ et $\PP(A\cup B)\ge\PP(B)$.

\emph{Cette borne est optimale.} Soient $a,b\in[0,1]$, disons $a\le b$. Sur $\Omega=\{1,2,3\}$, $\Tc=\Pc(\Omega)$, posons
$\PP(\{1\})=a$, $\PP(\{2\})=b-a$, $\PP(\{3\})=1-b$ (poids $\ge0$ de somme $1$, la $\sigma$-additivité est immédiate sur $\Omega$ fini : c'est une probabilité au sens de C1.10), $A=\{1\}$, $B=\{1,2\}$.
Alors $\PP(A)=a$, $\PP(B)=b$ et $\PP(A\cup B)=\PP(B)=b=\max(a,b)$ : aucune minoration meilleure, fonction de $\PP(A),\PP(B)$ seuls, n'est possible.
\begin{resultat}
$\max\bigl(\PP(A),\PP(B)\bigr)\le\PP(A\cup B)\ \bigl(\le\min(1,\PP(A)+\PP(B))\bigr)$, et la minoration est atteinte (cas $A\subset B$ ou $B\subset A$).
\end{resultat}
\textit{À ne pas confondre :} la minoration de l'\emph{intersection} $\PP(A\cap B)\ge\PP(A)+\PP(B)-1$ (C1.11, point 2), qui découle de $\PP(A\cup B)\le1$.
\end{exo}

% =====================================================================
\begin{exo}[Gérants bien ou mal informés]
\emph{Un gérant pris au hasard achète un titre ; le cours monte. Probabilité que le gérant soit mal informé ?}

\textbf{Modélisation.} On travaille sur un espace probabilisé $(\Omega,\F,\PP)$ avec les événements
$I=$ « le gérant choisi est bien informé » et $H=$ « le cours du titre monte ». Les données se traduisent par
(probabilités conditionnelles, C1.12) :
\[\PP(I)=\tfrac1{10}=0{,}1,\qquad \PP(H\mid I)=0{,}8,\qquad \PP(H^c\mid I^c)=0{,}6.\]
« Choisi au hasard » signifie que $\PP(I)$ est la proportion de gérants bien informés. On a $\PP(I)=0{,}1>0$ et $\PP(I^c)=0{,}9>0$
(C1.11, point 1), donc les conditionnements par $I$ et $I^c$ sont licites.

\begin{alerte}
L'énoncé donne la probabilité de \emph{baisse} pour un gérant mal informé. On suppose implicitement que le cours
\emph{monte ou baisse} (pas de cours inchangé), i.e.\ « baisse » $=H^c$. C'est l'interprétation retenue (et celle de la fiche TD1).
Si un cours stable était possible, les données ne suffiraient pas à déterminer $\PP(H\mid I^c)$.
\end{alerte}
% À VÉRIFIER : hypothèse implicite « baisse » = complémentaire de « hausse ».

\textbf{Étape 1 : $\PP(H\mid I^c)$.} $I^c=(H\cap I^c)\sqcup(H^c\cap I^c)$ (union disjointe), donc par additivité (C1.10 (iii)),
$\PP(H\cap I^c)+\PP(H^c\cap I^c)=\PP(I^c)$ ; en divisant par $\PP(I^c)>0$ (C1.12) :
\[\PP(H\mid I^c)=1-\PP(H^c\mid I^c)=1-0{,}6=0{,}4.\]

\textbf{Étape 2 : $\PP(H)$ (probabilités totales).} $\{I,I^c\}$ est une partition de $\Omega$ en événements de probabilité $>0$.
Comme $H=(H\cap I)\sqcup(H\cap I^c)$, l'additivité (C1.10 (iii)) puis la définition C1.12 ($\PP(H\cap I)=\PP(H\mid I)\PP(I)$) donnent
\[\PP(H)=\PP(H\mid I)\PP(I)+\PP(H\mid I^c)\PP(I^c)=0{,}8\times0{,}1+0{,}4\times0{,}9=0{,}08+0{,}36=0{,}44.\]
C'est la formule des probabilités totales (C1.13, \emph{Hors cours} : elle est redémontrée ici en une ligne).

\textbf{Étape 3 : Bayes.} Puisque $\PP(H)=0{,}44>0$, $\PP(I^c\mid H)$ est défini et, par C1.12 (deux fois),
\[\PP(I^c\mid H)=\frac{\PP(I^c\cap H)}{\PP(H)}=\frac{\PP(H\mid I^c)\PP(I^c)}{\PP(H)}=\frac{0{,}36}{0{,}44}=\frac{9}{11}.\]
(Formule de Bayes, C1.13 — \emph{Hors cours}, mais rien d'autre que la définition C1.12 n'est utilisé.)
Vérification numérique (awk) : $0{,}36/0{,}44=0{,}818182$.
\begin{resultat}
$\PP(\text{mal informé}\mid\text{le cours monte})=\PP(I^c\mid H)=\dfrac{9}{11}\approx0{,}818$.
\end{resultat}
\textit{Interprétation.} A priori $\PP(I^c)=0{,}9$ ; observer une hausse ne fait baisser cette probabilité que légèrement (à $0{,}818$),
car les mal informés sont très majoritaires et font tout de même monter le titre dans $40\,\%$ des cas.
Symétriquement, $\PP(I\mid H)=0{,}08/0{,}44=2/11\approx0{,}182$.
\end{exo}

\end{document}
